\chapter{Étape 4 : Implémentation et Tests}

\section{Structure du Projet}

\begin{verbatim}
grm-project/
├── backend/
│   ├── server.js              -- Point d'entrée Express
│   ├── package.json
│   ├── .env.example           -- Variables d'environnement
│   ├── database/
│   │   └── schema.sql         -- Schéma PostgreSQL + données initiales
│   └── src/
│       ├── config/
│       │   └── database.js    -- Pool PostgreSQL
│       ├── middleware/
│       │   └── auth.js        -- JWT authenticate + RBAC authorize
│       ├── controllers/       -- Logique métier (7 fichiers)
│       └── routes/            -- Routes REST (7 fichiers)
│
├── frontend/
│   ├── package.json
│   ├── public/index.html
│   └── src/
│       ├── App.js             -- Routing principal
│       ├── index.js           -- Point d'entrée React
│       ├── context/
│       │   └── AuthContext.js -- État global auth
│       ├── services/
│       │   └── api.js         -- Client Axios + intercepteurs
│       ├── components/
│       │   └── Layout/        -- Navigation + sidebar
│       └── pages/
│           ├── Login.js
│           ├── Dashboard.js
│           ├── Tenders/       -- Liste, Détail, Formulaire
│           ├── Resources/     -- Liste, Détail
│           ├── Suppliers/     -- Liste, Inscription
│           ├── Breakdowns/    -- Liste, Détail
│           └── Departments/   -- Liste et membres
│
└── docs/
    ├── main.tex
    ├── etape1_analyse.tex
    ├── etape2_conception.tex
    ├── etape3_architecture.tex
    └── etape4_implementation.tex
\end{verbatim}

\section{API REST — Endpoints Implémentés}

\subsection{Authentification}

\begin{table}[h!]
\centering
\begin{tabularx}{\textwidth}{|l|l|X|}
\hline
\textbf{Méthode} & \textbf{Route} & \textbf{Description} \\
\hline
POST & /api/auth/login & Authentification, retourne JWT \\
\hline
POST & /api/auth/register-supplier & Inscription fournisseur (utilisateur + société) \\
\hline
GET & /api/auth/profile & Profil utilisateur connecté \\
\hline
\end{tabularx}
\end{table}

\subsection{Appels d'Offre}

\begin{table}[h!]
\centering
\begin{tabularx}{\textwidth}{|l|l|X|}
\hline
\textbf{Méthode} & \textbf{Route} & \textbf{Description} \\
\hline
GET & /api/tenders & Liste tous les appels d'offre \\
\hline
GET & /api/tenders/:id & Détail avec besoins \\
\hline
POST & /api/tenders & Créer (responsable, chef) \\
\hline
PUT & /api/tenders/:id & Modifier (responsable) \\
\hline
PATCH & /api/tenders/:id/close & Fermer l'appel (responsable) \\
\hline
\end{tabularx}
\end{table}

\subsection{Offres}

\begin{table}[h!]
\centering
\begin{tabularx}{\textwidth}{|l|l|X|}
\hline
\textbf{Méthode} & \textbf{Route} & \textbf{Description} \\
\hline
GET & /api/offers/tender/:tenderId & Offres d'un appel (triées par prix) \\
\hline
POST & /api/offers/tender/:tenderId & Soumettre une offre (fournisseur) \\
\hline
PATCH & /api/offers/:id/accept & Accepter une offre (responsable) \\
\hline
PATCH & /api/offers/:id/reject & Rejeter une offre avec motif \\
\hline
\end{tabularx}
\end{table}

\subsection{Ressources}

\begin{table}[h!]
\centering
\begin{tabularx}{\textwidth}{|l|l|X|}
\hline
\textbf{Méthode} & \textbf{Route} & \textbf{Description} \\
\hline
GET & /api/resources & Liste toutes les ressources \\
\hline
GET & /api/resources/:id & Détail avec affectations et pannes \\
\hline
POST & /api/resources & Créer (responsable) \\
\hline
PUT & /api/resources/:id & Modifier (responsable) \\
\hline
DELETE & /api/resources/:id & Supprimer (responsable) \\
\hline
POST & /api/resources/:id/assign & Affecter à un département/personne \\
\hline
\end{tabularx}
\end{table}

\subsection{Pannes}

\begin{table}[h!]
\centering
\begin{tabularx}{\textwidth}{|l|l|X|}
\hline
\textbf{Méthode} & \textbf{Route} & \textbf{Description} \\
\hline
GET & /api/breakdowns & Liste toutes les pannes \\
\hline
GET & /api/breakdowns/:id & Détail avec constats \\
\hline
POST & /api/breakdowns & Signaler une panne (enseignant, chef) \\
\hline
POST & /api/breakdowns/:id/maintenance-report & Soumettre constat (technicien) \\
\hline
PATCH & /api/breakdowns/:id/return-to-supplier & Retourner au fournisseur \\
\hline
PATCH & /api/breakdowns/:id/resolve & Résoudre la panne \\
\hline
\end{tabularx}
\end{table}

\section{Extraits de Code}

\subsection{Middleware d'Authentification JWT}

\begin{lstlisting}[language=JavaScript, caption={src/middleware/auth.js}]
const authenticate = (req, res, next) => {
    const header = req.headers.authorization;
    if (!header) return res.status(401).json({ error: 'No token provided' });

    const token = header.split(' ')[1];
    try {
        req.user = jwt.verify(token, process.env.JWT_SECRET);
        next();
    } catch {
        res.status(401).json({ error: 'Invalid token' });
    }
};

const authorize = (...roles) => (req, res, next) => {
    if (!roles.includes(req.user.role))
        return res.status(403).json({ error: 'Access denied' });
    next();
};
\end{lstlisting}

\subsection{Transaction PostgreSQL — Acceptation d'une Offre}

\begin{lstlisting}[language=JavaScript, caption={src/controllers/offerController.js}]
const accept = async (req, res) => {
    const client = await pool.connect();
    try {
        await client.query('BEGIN');
        // Accepter l'offre selectionnee
        const { rows } = await client.query(
            `UPDATE offers SET status='accepted' WHERE id=$1 RETURNING *`,
            [req.params.id]
        );
        // Rejeter automatiquement les autres offres du meme appel
        await client.query(
            `UPDATE offers SET status='rejected'
             WHERE tender_id=$1 AND id != $2`,
            [rows[0].tender_id, req.params.id]
        );
        // Marquer l'appel d'offre comme attribue
        await client.query(
            `UPDATE tenders SET status='awarded' WHERE id=$1`,
            [rows[0].tender_id]
        );
        await client.query('COMMIT');
        res.json(rows[0]);
    } catch (err) {
        await client.query('ROLLBACK');
        res.status(500).json({ error: err.message });
    } finally {
        client.release();
    }
};
\end{lstlisting}

\subsection{Context React — Authentification}

\begin{lstlisting}[language=JavaScript, caption={src/context/AuthContext.js}]
const login = async (email, password) => {
    const { data } = await api.post('/auth/login', { email, password });
    localStorage.setItem('token', data.token);
    localStorage.setItem('user', JSON.stringify(data.user));
    api.defaults.headers.common['Authorization'] = `Bearer ${data.token}`;
    setUser(data.user);
    return data.user;
};
\end{lstlisting}

\section{Installation et Démarrage}

\subsection{Prérequis}

\begin{itemize}
    \item Node.js 18+
    \item PostgreSQL 14+
    \item npm ou yarn
\end{itemize}

\subsection{Configuration de la Base de Données}

\begin{lstlisting}[language=SQL, caption={Initialisation}]
-- Creer la base de donnees
CREATE DATABASE grm_db;

-- Executer le schema
psql -U postgres -d grm_db -f backend/database/schema.sql
\end{lstlisting}

\subsection{Démarrage du Backend}

\begin{lstlisting}[language=bash, caption={Backend setup}]
cd backend
cp .env.example .env
# Editer .env avec vos parametres PostgreSQL et JWT_SECRET

npm install
npm run dev
# Serveur demarre sur http://localhost:5000
\end{lstlisting}

\subsection{Démarrage du Frontend}

\begin{lstlisting}[language=bash, caption={Frontend setup}]
cd frontend
npm install
npm start
# Application demarre sur http://localhost:3000
\end{lstlisting}

\subsection{Comptes de Test}

\begin{table}[h!]
\centering
\begin{tabularx}{\textwidth}{|l|l|l|}
\hline
\textbf{Email} & \textbf{Mot de passe} & \textbf{Rôle} \\
\hline
responsable@faculte.ma & admin123 & Responsable des ressources \\
\hline
chef.info@faculte.ma & admin123 & Chef de département Informatique \\
\hline
ahmed.benali@faculte.ma & admin123 & Enseignant \\
\hline
tech@faculte.ma & admin123 & Technicien de maintenance \\
\hline
\end{tabularx}
\caption{Comptes prédéfinis pour les tests}
\end{table}

\section{Tests}

\subsection{Tests Unitaires — Middleware d'Authentification}

\begin{lstlisting}[language=JavaScript, caption={Tests JWT}]
describe('authenticate middleware', () => {
    it('should reject requests without token', () => {
        const req = { headers: {} };
        const res = { status: jest.fn().mockReturnThis(), json: jest.fn() };
        authenticate(req, res, jest.fn());
        expect(res.status).toHaveBeenCalledWith(401);
    });

    it('should accept valid JWT and set req.user', () => {
        const token = jwt.sign({ id: '1', role: 'responsable' }, process.env.JWT_SECRET);
        const req = { headers: { authorization: `Bearer ${token}` } };
        const next = jest.fn();
        authenticate(req, {}, next);
        expect(next).toHaveBeenCalled();
        expect(req.user.role).toBe('responsable');
    });
});
\end{lstlisting}

\subsection{Tests d'Intégration — API}

\begin{lstlisting}[language=JavaScript, caption={Tests API avec Supertest}]
describe('POST /api/auth/login', () => {
    it('should return JWT for valid credentials', async () => {
        const res = await request(app)
            .post('/api/auth/login')
            .send({ email: 'responsable@faculte.ma', password: 'admin123' });
        expect(res.status).toBe(200);
        expect(res.body).toHaveProperty('token');
        expect(res.body.user.role).toBe('responsable');
    });

    it('should return 401 for invalid credentials', async () => {
        const res = await request(app)
            .post('/api/auth/login')
            .send({ email: 'wrong@test.com', password: 'wrong' });
        expect(res.status).toBe(401);
    });
});
\end{lstlisting}

\section{Déploiement Docker}

\begin{lstlisting}[language=bash, caption={docker-compose.yml}]
version: '3.8'
services:
  db:
    image: postgres:15-alpine
    environment:
      POSTGRES_DB: grm_db
      POSTGRES_USER: postgres
      POSTGRES_PASSWORD: secret
    volumes:
      - ./backend/database/schema.sql:/docker-entrypoint-initdb.d/schema.sql
      - pgdata:/var/lib/postgresql/data

  backend:
    build: ./backend
    ports:
      - "5000:5000"
    environment:
      DB_HOST: db
      DB_NAME: grm_db
      DB_USER: postgres
      DB_PASSWORD: secret
      JWT_SECRET: change_this_in_production
    depends_on:
      - db

  frontend:
    build: ./frontend
    ports:
      - "80:80"
    depends_on:
      - backend

volumes:
  pgdata:
\end{lstlisting}

\begin{lstlisting}[language=bash, caption={Lancement complet}]
docker-compose up --build
# Application disponible sur http://localhost
\end{lstlisting}

\section{Conclusion}

Le système GRM implémente l'ensemble des fonctionnalités décrites dans le cahier des charges :

\begin{itemize}
    \item \checkmark Gestion complète des appels d'offre (CRUD + cycle de vie)
    \item \checkmark Soumission et sélection des offres fournisseurs avec notifications
    \item \checkmark Mise en liste noire des fournisseurs non conformes
    \item \checkmark Inventaire et affectation des ressources avec historique
    \item \checkmark Signalement et suivi des pannes
    \item \checkmark Constats de maintenance par les techniciens
    \item \checkmark Retour de ressources défectueuses aux fournisseurs
    \item \checkmark Authentification JWT et contrôle d'accès par rôle (RBAC)
    \item \checkmark Interface responsive accessible à tous les acteurs
\end{itemize}
